
\section{機能概要・前提条件・制約事項}

\subsection{機能概要}

本システムでは、大きく分けて競技者向け機能、主催者向け機能、観戦者向け機能を
提供します。競技者および主催者による大会運営を支援するとともに、
観戦者が特定の選手を中心として大会の進行状況や試合結果、過去の成績などを
確認できるようにすることでゲーム大会の観戦体験を支援します。

\subsubsection{競技者向け機能}

\paragraph{大会一覧表示・検索機能}
大会名、開催日時、開催場所、対象ゲーム、参加人数、参加受付状況などを一覧表示します。
大会名や対象ゲーム、開催日、参加受付状況等の条件による検索・絞り込みを行えます。

\paragraph{大会競技者登録機能}
競技者が大会を選択し、登録済みの競技者アカウントを使用して参加登録を行います。
初めて利用する場合はアカウントを作成し、競技者を識別するための情報と、
ゲーム内で使用するプレイヤー名などを登録します。競技参加登録情報はデータベースに保存します。

\paragraph{競技参加登録確認機能}
競技者自身が、参加登録した大会、登録状況、受付状況などを確認します。

\paragraph{大会進行状況確認機能}
大会全体の進行状況、試合中・試合待ち・試合終了などの状態を確認します。

\paragraph{対戦相手・試合予定確認機能}
競技者が自分の対戦相手、試合予定、対戦場所または試合番号などを確認します。

\paragraph{試合結果確認機能}
競技者が自分の試合結果や、トーナメント全体の結果を確認します。

\paragraph{競技者向け通知機能}
競技者がサイトを使わなくてもアカウントに登録したメールアドレスに試合予定の変更や試合結果の更新などメールで通知します。

\subsubsection{主催者向け機能}

\paragraph{大会作成機能}
大会名、対象ゲーム、開催日時、開催場所、最大参加人数、大会形式、参加受付期間などを登録します。
なお、大会形式はシングルエリミネーションまたは総当たりのどちらかとします。

\paragraph{競技者管理機能}
大会への競技者を一覧表示し、競技者名、登録状況、受付状況などを確認します。
また、競技者名や登録状況などの条件で参加者を検索・絞り込みを行います。

\paragraph{受付管理機能}
参加登録済みの競技者について、主催者が本人または登録情報を確認したうえで、管理画面から
「受付済み」に変更します。参加登録を行っただけでは受付完了にはならず、受付済み・未受付の状態を区別して表示します。

\paragraph{トーナメント作成機能}
参加登録された競技者数に応じた、シングルエリミネーションまたは総当たりのトーナメント表を作成します。
作成後は、対戦組み合わせや試合番号などを確認します。

\paragraph{試合結果登録機能}
試合終了後、競技者が自分の試合結果を入力し、対戦相手が入力内容を確認します。
両者の確認が一致した場合、試合結果を確定し、トーナメント表や選手の戦績に自動的に反映します。
結果が一致しない場合や、一定時間経過しても確認が行われない場合は、主催者が確認・修正できるようにします。
また、主催者が必要に応じて試合結果を直接登録できるようにします。

\paragraph{大会進行管理機能}
試合結果に基づいて優勝者や最終順位などを記録し、大会終了後も確認できるようにします。

\subsubsection{観戦者向け機能}

\paragraph{大会情報閲覧機能}
大会名、開催日時、対象ゲーム、競技者、現在の進行状況などを閲覧できます。
また、大会名、対象ゲーム、開催日時、開催状況などの条件で検索・絞り込み出来ます。

\paragraph{選手検索機能}
登録されている選手を、プレイヤー名などの条件で検索できます。
検索結果から選手のプロフィールや大会成績を確認できます。

\paragraph{お気に入り選手登録機能}
観戦者が興味のある選手をお気に入りに登録できます。
お気に入り情報は選手を識別するための固定IDと関連付けて管理し、
プレイヤー名が変更された場合でも同じ選手として追跡できるようにします。

\paragraph{選手追跡機能}
選手ごとのページにおいて、現在参加している大会、現在の試合、次の試合、今大会の成績、過去の大会成績などをまとめて確認できます。

\paragraph{試合状況確認機能}
大会全体の試合一覧から、試合中・試合待ち・試合終了などの状態や対戦相手、試合結果を確認できます。
また、お気に入り選手の試合を優先して確認できるようにします。

\paragraph{次回試合確認機能}
お気に入り選手について、次の対戦相手、試合予定、試合番号などを確認できます。
対戦相手が未確定の場合は、その状態を表示します。

\paragraph{選手成績確認機能}
選手ごとに、過去の大会への参加履歴、各大会の順位、試合結果などを確認できます。

\paragraph{対戦履歴確認機能}
選手を指定し、その選手と他の選手との過去の対戦結果を確認できます。

\paragraph{トーナメント状況確認機能}
大会全体のトーナメント表を表示し、各選手の勝ち上がり状況や次の対戦相手を確認できます。
観戦者はお気に入り選手だけでなく、他の選手の進行状況も確認できます。

\subsection{前提条件}
\begin{itemize}
  \item 主催者・競技者・観戦者がスマートフォンまたはPCを利用できることとします。
  \item システムを利用する端末がインターネットに接続できることとします。
  \item 大会の参加登録、受付、試合進行および結果確認に必要な情報が、主催者や参加者から提供されることとします。
  \item 選手の大会参加履歴や試合結果などを、選手を識別するための固定IDと関連付けてデータベースに保存・参照できることとします。
  \item 試合結果は、主催者または試合担当者などによる確認を経てシステムに反映されることとします。
\end{itemize}

\subsection{制約事項}
\begin{itemize}
  \item 収集する個人情報はシステムの運用に必要な範囲に限定し、公開用の選手情報と非公開の個人情報を区別します。
  \item 大会情報や試合結果の登録・変更は主催者などの権限を持つ利用者に限定し、一般の参加者や観戦者は閲覧のみ可能とします。
  \item ログイン情報を安全に管理し、パスワードを平文で保存しないこととします。
  \item 通信障害やシステム障害に備え、主催者が試合結果を確認・再登録できるようにすることとします。
  \item 開発期間を考慮し、対応する大会形式はトーナメント方式を中心とした限定的なものとすることとします。
  \item 大会規模は、開発期間及びシステムの処理能力を考慮して設定することとします。
  \item 試合の映像配信機能は実装対象とせず、試合状況、トーナメント表、結果などの情報提供を中心とすることとします。
  \item 観戦者向け機能は、お気に入り選手の追跡、現在・次の試合の確認、過去成績及び対戦履歴の閲覧を優先します。プッシュ通知などの追加機能は開発状況に応じて検討することとします。
\end{itemize}
